
D\'
emontrons que (i) implique (ii), (iii), et (iv). Supposons donc que @@M0@@ soit limit\'ee par un faisceau inversible sur un @@M1@@. Alors @@M2@@ satisfait \`a (iv) par EGA II 4.6.12. Supposons @@M3@@ plat sur @@M4@@. Alors le polyn\^ome en @@M5@@ et @@M6@@, @@M7@@ ne d\'epend que de la composante connexe de @@M8@@ (EGA III 7.9.4). Par suite @@M9@@ satisfait \`a (ii), et compte tenu de 2.1, \`a (iii).

Notons ensuite que la famille des classes des @@M10@@ (resp. @@M11@@) est limit\'ee, donc satisfait \`a (iii) (resp. \`a (ii)).

Supposons que @@M12@@ satisfasse \`a (iv). Alors @@M13@@, d'o\`u @@M14@@, et de m\^eme, @@M15@@. De plus, @@M16@@, d'o\`u @@M17@@. Donc @@M18@@ satisfait \`a (iii).

D\'emontrons enfin que (ii) (resp. (iii)) implique (i). Gr\^ace \`a 2.10 et 2.3, on peut remplacer @@M19@@ par un multiple\,; supposons donc @@M20@@ tr\`es ample. \'Evidemment on peut couper @@M21@@ en morceaux\,; supposons donc qu'il existe une section lin\'eaire @@M22@@ de @@M23@@, qui est int\`egre sur @@M24@@, \`a fibres de dimension @@M25@@. En vertu de 3.8, on peut remplacer @@M26@@ par @@M27@@, (compte tenu de 1.7 (resp. la formule de projection), qui implique que la condition (ii) (resp. (iii)) sur @@M28@@ est \'equivalente \`a (ii) (resp. \`a (iii)) sur @@M29@@)\,; supposons donc @@M30@@. Si @@M31@@ est @@M32@@, on peut remplacer @@M33@@ par la famille des classes des @@M34@@ gr\^ace \`a 2.3 et 2.10\,; supposons donc @@M35@@.

Si dans les polyn\^omes de Hilbert @@M36@@, @@M37@@ reste born\'e et @@M38@@ reste minor\'e, alors les trois coefficients restent born\'es gr\^ace \`a 2.3 et 2.7\,; donc @@M39@@ est limit\'ee en vertu de 2.11\,; ainsi (ii) implique
